\chapter*{Conclusion générale}
\addcontentsline{toc}{chapter}{Conclusion générale}
\markboth{Conclusion générale}{}

\section*{Bilan}
Ce projet de fin d'études a porté sur la conception et la mise en œuvre d'une solution décisionnelle complète au service du CETUD, visant à exploiter l'Enquête Ménages-Déplacements et les données de comptage de trafic de Dakar pour éclairer la prise de décision en matière de mobilité urbaine. La démarche suivie s'est appuyée sur la méthodologie GIMSI, qui a structuré le travail en quatre grandes étapes~: l'identification des besoins et des acteurs (chapitre~\ref{ch:existant}), la conception de la solution et du modèle décisionnel (chapitres~\ref{ch:conception} et~\ref{ch:decisionnel}), la mise en œuvre technique de la plateforme et des modèles prédictifs (chapitres~\ref{ch:dev} et~\ref{ch:ml}), et l'évaluation des résultats obtenus (chapitre~\ref{ch:resultats}).

\section*{Contributions}
Au-delà du volet décisionnel classique (ETL Talend, entrepôt de données en schéma constellation, cockpit Knowage BI), le projet a été enrichi par le développement de quatre modèles de Machine Learning~: la segmentation des usagers par K-Means, la recommandation de mode de transport par Random Forest, la détection d'anomalies de trafic par un consensus de trois algorithmes (Z-Score, Isolation Forest, Local Outlier Factor) complété par un clustering spatial DBSCAN, et la prédiction du risque d'inaccessibilité par Gradient Boosting. Ces modèles ont été exposés au sein d'une application web complète, combinant un backend FastAPI sécurisé par authentification JWT et un frontend React, à destination à la fois des citoyens et des décideurs du CETUD (chef d'exploitation et directeur de planification). Sur le plan des compétences acquises, ce projet a permis une mise en pratique transversale couvrant l'ensemble de la chaîne de valeur de la donnée~: de la modélisation décisionnelle classique jusqu'aux techniques avancées de Machine Learning, en passant par le développement complet d'une application web sécurisée. Cette double dimension, décisionnelle et prédictive, constitue l'apport principal du projet par rapport à une solution de reporting classique.

\section*{Difficultés rencontrées}
La conduite de ce projet a confronté à plusieurs difficultés caractéristiques d'un projet décisionnel de cette ampleur~: tout d'abord, la \textbf{confidentialité partielle des données du CETUD}, qui n'a pas pu mettre à disposition l'ensemble des sources souhaitées, contraignant certains indicateurs à être estimés ou partiellement renseignés~; ensuite, la \textbf{complexité structurelle et documentaire des données} issues de l'EMD (codes de variables non auto-documentés, hétérogénéité des formats, incomplétude de certaines variables), nécessitant un travail conséquent de nettoyage et de documentation avant toute exploitation~; l'arbitrage entre granularité fine des données individuelles et performance des agrégations dans l'entrepôt, qui a motivé le choix du schéma en constellation~; le déséquilibre des classes dans la tâche de prédiction d'inaccessibilité, traité par sur-échantillonnage SMOTE~; et la nécessité de concilier deux chaînes de traitement distinctes sans rompre la cohérence d'ensemble de la solution.

\section*{Perspectives d'amélioration}
Plusieurs axes d'amélioration et d'extension sont envisageables pour la suite~: le développement d'une \textbf{application mobile} (iOS / Android) permettant aux citoyens d'accéder aux recommandations de transport en situation de mobilité~; l'\textbf{extension géographique} de la solution à d'autres agglomérations sénégalaises ou sous-sahariennes, en réutilisant le pipeline ETL et le modèle décisionnel conçus~; l'\textbf{intégration de données de trafic en temps réel} (issues des opérateurs de transport ou de capteurs IoT) pour permettre une actualisation continue de l'entrepôt et des tableaux de bord, au lieu d'un chargement batch périodique~; la mise en place de \textbf{rapports mensuels automatisés} (génération PDF programmée, envoi par email aux décideurs du CETUD)~; et enfin, l'exploration de techniques de \textbf{Deep Learning}, notamment les réseaux de neurones récurrents de type LSTM pour la prédiction temporelle du trafic, en remplacement ou en complément du modèle d'anomalies actuel.

En définitive, ce projet démontre la faisabilité et l'intérêt d'une approche combinant Business Intelligence traditionnelle et Machine Learning pour répondre aux enjeux de mobilité urbaine d'une métropole comme Dakar, et pose les bases d'un système qui pourra être consolidé et étendu dans le cadre d'une collaboration future avec le CETUD.

\clearpage
